% =====================================================================
%  Tóm tắt bài báo cáo (TV1)
%  Chỉ nêu kết quả đã kiểm chứng. Số BER ở đoạn cuối là giá trị LÝ THUYẾT
%  (tính chính xác bằng vét cạn 128 mẫu lỗi, đối chiếu mô phỏng 400.000
%  khối): BER = tỉ lệ bit DỮ LIỆU sai sau giải mã. Khi Chương 5 (TV9) có
%  bảng mô phỏng, kiểm tra các số trong Bảng 5.3 khớp với đoạn này (cùng
%  định nghĩa BER, sai khác chỉ do nhiễu mô phỏng).
% =====================================================================
\chuongphu{TÓM TẮT BÀI BÁO CÁO}

Báo cáo trình bày mã Hamming (7,4) dưới góc nhìn đại số tuyến tính trên trường
$\GF$ và xây dựng phần mềm mô phỏng quá trình mã hóa, gây nhiễu, phát hiện và
sửa lỗi cho chuỗi ký tự. Mã được xét ở dạng hệ thống: khối dữ liệu $4$ bit $u$
được mã hóa thành từ mã $7$ bit $c = uG$ với ma trận tạo mã
$G = [\,I_4 \mid P\,]$; bên nhận tính hội chứng $s = r\HT$ với ma trận kiểm tra
$H = [\,\PT \mid I_3\,]$ để xác định và sửa một bit sai.

Phần lý thuyết chứng minh $G\HT = 0$ và khoảng cách Hamming tối thiểu
$\dmin = 3$ từ các cột của $H$, rồi giải thích vì sao mã (7,4) là mã hoàn hảo:
sửa được mọi mẫu một lỗi nhưng sửa nhầm mọi mẫu hai lỗi, nên không thể vừa sửa
một lỗi vừa phát hiện hai lỗi. Để khắc phục, báo cáo xét mã Hamming mở rộng
(8,4) có thêm bit kiểm tra tổng thể, cho phép phát hiện lỗi kép. Các ví dụ tính
tay được thực hiện trên khối dữ liệu $1011$ và trên ký tự \code{'A'}.

Phần mềm viết bằng Python gồm các module mã hóa và giải mã một khối theo mã
(7,4) và mã mở rộng (8,4), chuyển đổi
giữa chuỗi ký tự và các khối $4$ bit, giả lập kênh đối xứng nhị phân (BSC) và
kiểm thử tự động. Kiểm tra vét cạn trên $16$ khối dữ liệu và $128$ mẫu lỗi trong
một từ mã cho thấy bộ giải mã (7,4) khôi phục đúng dữ liệu khi và chỉ khi có
không quá một bit lỗi. Trên kênh BSC, tỉ lệ lỗi bit dữ liệu sau giải mã (BER)
tính theo lý thuyết giảm từ $0{,}01$ xuống khoảng $0{,}00087$ khi xác suất đảo
bit $p = 0{,}01$, còn khoảng $0{,}0194$ khi $p = 0{,}05$ và $0{,}0669$ khi
$p = 0{,}1$; lợi ích của mã giảm dần khi kênh xấu đi.

\vspace{0.5em}
\noindent\textbf{Từ khóa:} mã Hamming (7,4), mã Hamming mở rộng (8,4), trường
$\GF$, mã khối tuyến tính, ma trận tạo mã, ma trận kiểm tra, hội chứng, khoảng
cách Hamming, kênh đối xứng nhị phân.
